خلال الربع محل الدراسة (أكتوبر - ديسمبر 2025)، قادت \textbf{الشركات الكبيرة} التعافي في مؤشر بارومتر الأعمال؛ حيث تجاوز مؤشر الأعمال \textbf{لها المستوى المحايد} بـ 16 نقطة، مسجلا قيما أفضل من \textbf{الربعين السابق والمناظر}، ويعكس ذلك تحسنا ملحوظا في المؤشرات الفرعية، وتحديدا \textbf{الإنتاج، والمبيعات المحلية، والصادرات، ومستوى استغلال الطاقة الإنتاجية}. وفي المقابل، جاء التعافي في \textbf{الشركات الصغيرة والمتوسطة} بوتيرة أبطأ من الشركات الكبيرة، وتجاوز مؤشر البارومتر \textbf{لهذه الشركات المستوى المحايد} بـ 4 نقاط وتخطيه الربع المناظر بـ 6 نقاط، إلا أنه جاء أقل من الربع السابق، مما يعكس ارتفاع مؤشرات \textbf{الإنتاج، والمبيعات المحلية}، \textbf{في حين لم تحقق الصادرات واستغلال الطاقة الإنتاجية} نفس مستوى الأداء (الشكل 3-1).
